% IEEE conference paper: Field Consistency: The Cost of Enforcing PDE Residuals on 2D Fields
% Compile: pdflatex ieee_paper.tex   (figures in ../figs)
\documentclass[conference,10pt]{IEEEtran}

\usepackage{graphicx}
\usepackage{amsmath,amssymb}
\usepackage{booktabs}
\usepackage[table]{xcolor}
\usepackage{url}
\usepackage[margin=1in]{geometry}

\begin{document}

\title{Field Consistency: The Cost of Enforcing PDE Residuals on 2D Fields}

\author{\IEEEauthorblockN{Sehaj Randhir Singh}
\IEEEauthorblockA{\textit{Department of Electrical and Computer Engineering,}\\
\textit{NYU Tandon School of Engineering (independent researcher)}\\
Brooklyn, NY, USA\\
sehajrsinghs@gmail.com}}

\maketitle

\begin{abstract}
The tensor pipeline extends to 2D fields: the same transformer with a field-shaped head learns the full $u(x,y)$ surface of a heat plate. The showcase number -- 5.9\% peak field error -- is real and is the easiest member of the validation distribution, whose honest mean is ~29\%; this paper reports the distribution, not just the showcase. Enforcing the Poisson residual cuts it 6--8$\times$ while field fidelity degrades: consistency is not accuracy at the field level. Against a per-instance PINN (DeepXDE, ~19 minutes per problem), the specialist wins on its own instance at 0.03--0.06\% error but with a higher governing-equation residual; the generalist answers any instance in one forward pass. Both statements are true; the claim is the tradeoff.
\end{abstract}

\begin{IEEEkeywords}
physics-informed neural networks; PDEs; PINN; fields; verification
\end{IEEEkeywords}

\section{Introduction}
At the field level the pattern from trajectory domains repeats and sharpens: enforcing the governing equation makes predictions consistent, not correct. The two must be measured separately, and here both are -- including the honest distribution number that a showcase-only evaluation hides.

\section{Protocol}
Heat-plate domain ($u_{xx}+u_{yy}=f$), field-shaped head, 240 epochs, canonical showcase plus the full validation distribution; Burgers' equation for the per-instance DeepXDE comparison (per-instance PINN at 0.03--0.06\% full-field error vs. one-pass generalist at 33--51\% with lower governing-equation residual).

\section{Results}
\begin{table}[!t]
\caption{Field-level generalist vs.\ per-instance PINN (DeepXDE). A
specialist wins on its own problem; the generalist is a consistent
solver of all instances.}
\label{tab:field}
\centering
\begin{tabular}{lcc}
\toprule
Instance & DeepXDE & generalist \\
\midrule
instances & -- & -- \\
\bottomrule
\end{tabular}
\end{table}

\begin{thebibliography}{1}
\bibitem{lu2021deepxde} L.~Lu, X.~Meng, Z.~Mao, and G.~E. Karniadakis, ``DeepXDE: a deep learning library for solving differential equations,'' SIAM Review, vol.~63, pp.~208--228, 2021.
\end{thebibliography}

\end{document}
